% GENERATED FILE. Edit canonical lessons, then run npm run generate:books.
% canonical-source-hash: fnv1a64:76ab14bb8e44a37d
% canonical-lessons: ES-C445-jubilarse, ES-C445-responsabilidad, ES-C445-exigente, ES-C445-ensenar, ES-C445-repetir, ES-R445-repaso-puesto, ES-C445-sintesis-puesto

\chapter{El Puesto}
\label{ch:es-puesto}
\hlchaptermodality{445}

\begin{chapteropening}
\textbf{By the end of this chapter:} \emph{I can read a job advert and say what the work asks of me: what it demands, what I am answerable for, what I would be teaching, and when the person I am replacing leaves.}

Read the message or notice this chapter's words exist for, then say the same thing back.
\end{chapteropening}

\section[jubilarse]{jubilarse — the shout at the end of the work}

\label{lesson:ES-C445-jubilarse}



\begin{quote}
Say \textbf{jubilado} and \textbf{el trabajo}. You have the word for the person who has stopped. Name the act of stopping.
\end{quote}

\subsection*{You'll want to know: jubilarse}
\textbf{jubilarse} — to retire.

\emph{Se jubila en marzo} — he retires in March. \emph{Me jubilé hace dos años} — I retired two years ago. \textbf{la jubilación} is both the act and the pension.

It is reflexive: you retire \emph{yourself}. \emph{Jubilar} without the pronoun is what an employer does to somebody, and it can sound abrupt.

\begin{cousinweb}
Latin \textbf{iubilāre} — \textbf{to shout, to whoop for joy}. It was a countryman's word for the noise you make across a field, not a word about work at all.

The connection runs through the \textbf{jubilee}: in the older calendar a jubilee year cancelled debts and released people from service, and \emph{jubilar} came to mean \emph{to release from service}. The shout came first and the pension came after.

So \textbf{el jubilado} is, underneath, \emph{the one who has been cheered off}, and every Spanish retirement notice carries a faint sound of celebration nobody hears any more.

English keeps the shout in \textbf{jubilant} and \textbf{jubilation}, and has no verb for the Spanish sense at all.
\end{cousinweb}

\subsection*{Your turn}
Say these aloud:

\begin{itemize}
\raggedright
  \item \emph{jubilarse}
  \item she retires in March after forty years of work
  \item what \emph{iubilāre} meant
\end{itemize}

\begin{tcolorbox}[breakable,colback=teal!4,colframe=teal!35!black,title={Before you move on}]
To retire? (\textbf{Jubilarse}.) Is it reflexive? (\textbf{Yes}.) What did the Latin verb mean? (\textbf{To shout for joy}.)
\end{tcolorbox}
\section[la responsabilidad]{la responsabilidad — answerable}

\label{lesson:ES-C445-responsabilidad}



\begin{quote}
Say \textbf{el contrato} and \textbf{jubilarse}. A contract lists what a job pays and what it asks. Name what it asks.
\end{quote}

\subsection*{You'll want to know: la responsabilidad}
\textbf{la responsabilidad} — responsibility.

\emph{Es una responsabilidad grande} — it is a big responsibility. \emph{Bajo mi responsabilidad} — on my own responsibility. The adjective is \textbf{responsable}, which is both \emph{responsible} and, in an advert, \emph{the person in charge}: \emph{el responsable de compras}.

Six syllables, stress on the last: \emph{res-pon-sa-bi-li-DAD}.

\begin{cousinweb}
\textbf{re-} back, plus \textbf{spondēre}, to promise, plus the ending \textbf{-bilidad} that turns an adjective into a quality.

To be \emph{responsable} is to be \textbf{able to promise back} — to answer for what you did. That is the same \emph{spondēre} you already have in a place you would never look for it:

\noindent
\begin{tabularx}{\linewidth}{@{}>{\raggedright\arraybackslash}X>{\raggedright\arraybackslash}X@{}}
\toprule
\textbf{word} & \textbf{the promise in it} \\
\midrule
\textbf{el esposo} & the one who was \textbf{promised} \\
responder & to promise \textbf{back}, to answer \\
la responsabilidad & being answerable \\
\bottomrule
\end{tabularx}

A spouse and a responsibility are the same verb. Latin \emph{spondēre} was the formal word for pledging — a betrothal was a \emph{sponsalia} — and English keeps it in \textbf{sponsor}, \textbf{respond} and \textbf{despondent}, which is what you are when the promise has gone out of you.
\end{cousinweb}

\begin{culture}
The long ending is worth a moment, because it is a machine.

\textbf{-idad} (and \textbf{-bilidad} after \emph{-ble}) turns an adjective into the abstract noun of its quality, and it is completely regular:

\noindent
\begin{tabularx}{\linewidth}{@{}>{\raggedright\arraybackslash}X>{\raggedright\arraybackslash}X@{}}
\toprule
\textbf{adjective} & \textbf{quality} \\
\midrule
responsable & responsabilidad \\
posible & posibilidad \\
fácil & facilidad \\
\bottomrule
\end{tabularx}

English does the same with \emph{-ity}, and the two endings are the same Latin \emph{-itātem}, which is why they line up so neatly. If you know a Spanish adjective in \textbf{-ble}, you can usually build its noun without being taught it.

Note the stress: every one of these nouns ends stressed, on the \textbf{-dad}.
\end{culture}

\subsection*{Your turn}
Say these aloud:

\begin{itemize}
\raggedright
  \item \emph{la responsabilidad}
  \item it is a big responsibility and he retires in March
  \item why \emph{esposo} and \emph{responsabilidad} share a root
\end{itemize}

\begin{tcolorbox}[breakable,colback=teal!4,colframe=teal!35!black,title={Before you move on}]
Responsibility? (\textbf{La responsabilidad}.) What does \emph{spondēre} mean? (\textbf{To promise}.) Which syllable is stressed? (\textbf{The last}.)
\end{tcolorbox}
\section[exigente]{exigente — driving it out of you}

\label{lesson:ES-C445-exigente}



\begin{quote}
Say \textbf{difícil} and \textbf{fácil}. Say \textbf{la responsabilidad}. A job can be hard without anybody being hard on you. Name the word for a boss who is.
\end{quote}

\subsection*{You'll want to know: exigente}
\textbf{exigente} — demanding, exacting.

\emph{Un jefe exigente} — a demanding boss. \emph{Es muy exigente consigo mismo} — he is very hard on himself. The verb is \textbf{exigir}, to demand: \emph{exigen dos años de experiencia} — they require two years of experience, which is the line you meet in a job advert.

It does not change for gender: \emph{un jefe exigente}, \emph{una jefa exigente}.

\begin{cousinweb}
\textbf{ex-} out, plus \textbf{-igente} from Latin \textbf{agere}, to drive.

\emph{Exigere} was \textbf{to drive something out} — originally to drive out a payment, to collect what was owed. An \emph{exigente} boss drives the work out of you, and the picture is the same one the Roman tax collector used.

\emph{Agere} is one of the great Latin roots, and the words it gave English all keep some sense of driving something forward:

\noindent
\begin{tabularx}{\linewidth}{@{}>{\raggedright\arraybackslash}X>{\raggedright\arraybackslash}X@{}}
\toprule
\textbf{word} & \textbf{what is driven} \\
\midrule
\textbf{agenda} & the things to be driven through \\
\textbf{actor} & the one who drives the action \\
\textbf{agile} & easily driven \\
\bottomrule
\end{tabularx}
\end{cousinweb}

\begin{culture}
\textbf{Difícil} and \textbf{exigente} are not the same complaint, and a job advert uses them differently.

\emph{Difícil} is a property of the task: the work is hard. \emph{Exigente} is a property of whoever is asking: the standard is high. A puzzle is \emph{difícil}; a person is \emph{exigente}. You can have an easy job with a demanding manager, and Spanish requires you to pick which you mean.

The ending \textbf{-ente} is doing that work. It comes from a Latin present participle and marks \textbf{the one doing the verb} — \emph{exigente} is the one demanding, the way \emph{estudiante} is the one studying and \emph{cliente} is the one leaning on you for service.

So when you meet an unfamiliar word in \textbf{-ante} or \textbf{-ente}, look for a verb inside it and read the whole as \emph{the one who does that}.
\end{culture}

\subsection*{Your turn}
Say these aloud:

\begin{itemize}
\raggedright
  \item \emph{exigente}
  \item an easy job with a demanding boss
  \item what the two pieces of \emph{exigir} mean
\end{itemize}

\begin{tcolorbox}[breakable,colback=teal!4,colframe=teal!35!black,title={Before you move on}]
Demanding? (\textbf{Exigente}.) The verb? (\textbf{Exigir}.) What does \emph{agere} mean? (\textbf{To drive}.)
\end{tcolorbox}
\section[enseñar]{enseñar — to put a sign into somebody}

\label{lesson:ES-C445-ensenar}



\begin{quote}
Say \textbf{aprender} and \textbf{la clase}. Say \textbf{exigente}. You have the half of the class that receives. Name the half that gives.
\end{quote}

\subsection*{You'll want to know: enseñar}
\textbf{enseñar} — to teach, and equally \textbf{to show}.

\emph{Enseña matemáticas} — she teaches maths. \emph{¿Me enseñas la foto?} — will you show me the photo? \emph{Me enseñó a conducir} — he taught me to drive, with \textbf{a} before the second verb.

One word for both jobs, where English keeps \emph{teach} and \emph{show} apart.

\begin{cousinweb}
\textbf{en-} in, plus \textbf{señ-} from Latin \textbf{signum}, a sign or mark.

To \emph{enseñar} is \textbf{to mark something into somebody} — to put a sign in them. The Latin was \emph{insignāre}, and once you see it, \emph{teach} and \emph{show} stop being two meanings. Both are \textbf{putting a mark where somebody can see it}: on a wall, or in a head.

\noindent
\begin{tabularx}{\linewidth}{@{}>{\raggedright\arraybackslash}X>{\raggedright\arraybackslash}X@{}}
\toprule
\textbf{word} & \textbf{the sign in it} \\
\midrule
la señal & the sign itself \\
enseñar & to put the sign in \\
diseñar & to mark it \textbf{out} — to design \\
\bottomrule
\end{tabularx}

English took the same \emph{signum} through French into \textbf{ensign}, a banner, and into \textbf{design}, which is \emph{de-} plus the mark.
\end{cousinweb}

\begin{culture}
That one verb covering two English ones is not vagueness; it is a different cut of the same territory, and it goes the other way too.

Spanish splits what English joins all the time — \emph{ser} and \emph{estar} for one \emph{to be}, \emph{saber} and \emph{conocer} for one \emph{to know}. Here it joins what English splits. Neither language is being imprecise; they have drawn the lines in different places, and the only reliable way through is to learn where the line is rather than to translate across it.

A practical consequence: if you want \emph{show} specifically and fear being misheard as \emph{teach}, Spanish has \textbf{mostrar} — but in speech \emph{enseñar} is far commoner for both, and the context does the work.
\end{culture}

\subsection*{Your turn}
Say these aloud:

\begin{itemize}
\raggedright
  \item \emph{enseñar}
  \item she teaches maths, and I am learning
  \item will you show me the photo
\end{itemize}

\begin{tcolorbox}[breakable,colback=teal!4,colframe=teal!35!black,title={Before you move on}]
To teach? (\textbf{Enseñar}.) And to show? (\textbf{Enseñar} — the same word.) What does \emph{signum} mean? (\textbf{A sign, a mark}.)
\end{tcolorbox}
\section[repetir]{repetir — seeking it again}

\label{lesson:ES-C445-repetir}



\begin{quote}
Say \textbf{otra vez}. Say \textbf{enseñar}. You did not catch what the teacher said. Name what you ask for.
\end{quote}

\subsection*{You'll want to know: repetir}
\textbf{repetir} — to repeat.

\emph{¿Puede repetir, por favor?} — could you repeat that, please. This is the single most useful sentence in a language you are still learning, and it is worth saying aloud until it arrives without thinking.

It breaks its stem: \textbf{repito}, \textbf{repites}, \textbf{repite}, but \emph{repetimos} — the \textbf{e$\to$i} change, the one \emph{pedir} makes too.

\begin{cousinweb}
\textbf{re-} again, plus \textbf{petere}, to seek or to ask for.

And you have met \emph{petere} before, under a word from a completely different part of life. \textbf{El pedido} — the order you place with a shop — is built on \emph{pedir}, and \emph{pedir} is \emph{petere}.

\noindent
\begin{tabularx}{\linewidth}{@{}>{\raggedright\arraybackslash}X>{\raggedright\arraybackslash}X@{}}
\toprule
\textbf{word} & \textbf{the seeking in it} \\
\midrule
pedir, el pedido & to ask for a thing \\
\textbf{repetir} & to seek it \textbf{again} \\
\bottomrule
\end{tabularx}

That is why the two verbs break their stems the same way: \emph{pido}, \emph{pides}, \emph{pide}; \emph{repito}, \emph{repites}, \emph{repite}. They are not merely similar. They are one Latin verb, with and without a prefix, and the shared conjugation is the evidence rather than a coincidence you have to remember.
\end{cousinweb}

\begin{culture}
Notice what this chapter has quietly done. \textbf{Responsabilidad} shares a root with \textbf{el esposo}, a word from family; \textbf{repetir} shares one with \textbf{el pedido}, a word from shopping. Neither link is visible in modern Spanish, and neither is decoration — both times the shared root predicts something usable: the promise inside \emph{responsable}, and the conjugation of \emph{repetir}.

That is the return on etymology. Not that the stories are charming, but that a root you own tells you how an unfamiliar word will behave before anybody explains it.
\end{culture}

\subsection*{Your turn}
Say these aloud:

\begin{itemize}
\raggedright
  \item \emph{¿Puede repetir, por favor?}
  \item the teacher repeated it and then I understood
  \item \emph{pido} and \emph{repito}, and why they change the same way
\end{itemize}

\begin{tcolorbox}[breakable,colback=teal!4,colframe=teal!35!black,title={Before you move on}]
To repeat? (\textbf{Repetir}.) The \emph{yo} form? (\textbf{Repito}.) Which word you already had shares its root? (\emph{\textbf{El pedido}}, through \emph{pedir}.)
\end{tcolorbox}
\section[(jubilarse, responsabilidad, exigente …]{(jubilarse, responsabilidad, exigente, enseñar, repetir) — from cold}

\label{lesson:ES-R445-repaso-puesto}



\begin{quote}
Close the lessons before this one. Say all five: \emph{to retire}, \emph{responsibility}, \emph{demanding}, \emph{to teach}, \emph{to repeat}.
\end{quote}

\begin{grammarlens}[title={Grammar Lens: a root you own predicts how a new word behaves}]
Two of these words turned out to be built on roots you had already, from somewhere else entirely.

\noindent
\begin{tabularx}{\linewidth}{@{}>{\raggedright\arraybackslash}X>{\raggedright\arraybackslash}X>{\raggedright\arraybackslash}X@{}}
\toprule
\textbf{new word} & \textbf{the word you already had} & \textbf{what the shared root buys} \\
\midrule
responsabilidad & \textbf{el esposo} & the promise inside \emph{responsable} \\
repetir & \textbf{el pedido} & the \textbf{e$\to$i} stem change \\
\bottomrule
\end{tabularx}

\emph{Repetir} is the sharper of the two. \emph{Pedir} changes \emph{pido}, \emph{pides}, \emph{pide}; \emph{repetir} changes \emph{repito}, \emph{repites}, \emph{repite}. Nobody has to tell you that, because they are one Latin verb with a prefix on one of them.

That is the return on etymology here — not that the stories are pleasant, but that a root you own tells you how an unfamiliar word will behave before anybody explains it.
\end{grammarlens}

\subsection*{Your turn}
Say these aloud:

\begin{itemize}
\raggedright
  \item could you repeat that, please
  \item she teaches maths and she is very demanding
  \item he retires in March
\end{itemize}

\begin{grammarlens}[title={What you've built}]
\textbf{You can now read a job advert and say what the work asks of you} — what it demands, what you are answerable for, what you would be teaching or shown, and when the person you are replacing is leaving.

You also have two endings that keep paying. \textbf{-idad} turns an adjective into its quality and is regular, so \emph{posible} gives \emph{posibilidad}. And \textbf{-ente} or \textbf{-ante} marks \emph{the one doing the verb}, so an unfamiliar word in that shape has a verb inside it worth looking for.
\end{grammarlens}

\begin{tcolorbox}[breakable,colback=teal!4,colframe=teal!35!black,title={Before you move on}]
Responsibility? (\textbf{La responsabilidad}.) To teach, and to show? (\textbf{Enseñar}, for both.) To retire? (\textbf{Jubilarse}.)
\end{tcolorbox}
\section[Practice]{(una oferta y una despedida) — read it, then do it yourself}

\label{lesson:ES-C445-sintesis-puesto}



\begin{quote}
Say \emph{la responsabilidad} and \emph{jubilarse}. The two texts below are about the same desk: one person is leaving it and somebody else is wanted for it.
\end{quote}

\subsection*{Your turn}
Read both, then answer.

\textbf{1 — oferta de empleo}

\begin{quote}
Se busca persona \textbf{responsable} para el turno de tarde. Se \textbf{exige} experiencia previa y se valorará saber \textbf{enseñar} a personal nuevo. Contrato fijo. Interesados, enviar currículum antes del día 30.
\end{quote}

\textbf{2 — una nota en el tablón}

\begin{quote}
Después de treinta y un años, me \textbf{jubilo} el viernes. Ha sido un placer. Dejo el turno en buenas manos. Os espero a todos a las seis en el bar de enfrente.
\end{quote}

Say these aloud:

\begin{itemize}
\raggedright
  \item what the job requires, and what it merely prefers
  \item how long the second person has worked there
  \item which of the two texts has a person in it
\end{itemize}

\begin{culture}
The advert separates two things English often blurs.

\textbf{Se exige} experiencia — it is \emph{required}, and without it you are not eligible. \textbf{Se valorará} saber enseñar — it will be \emph{valued}, a point in your favour and nothing more. A reader who treats those as the same sentence either rules themselves out for no reason or applies for something they cannot do.

That distinction is the single most useful thing on the page, and it is carried entirely by the verb.

Notice too that the advert has \textbf{no people}: \emph{se busca}, \emph{se exige}, \emph{se valorará}, and even \emph{interesados} is a plural adjective standing in for a noun. The leaving note is the opposite — \emph{me jubilo}, \emph{dejo}, \emph{os espero}, first person throughout. One desk, two registers, and you can now read both.
\end{culture}

\begin{tcolorbox}[breakable,colback=teal!4,colframe=teal!35!black,title={Before you move on}]
Now your turn, out loud and then in writing. Five sentences applying for that job: say you are responsible, say you have the experience they require, say you have taught new staff before, ask them to repeat the closing date, and say you are sorry the other person is retiring.

Then check yourself: did you keep \textbf{se exige} and \textbf{se valorará} apart?
\end{tcolorbox}
